% !TEX program = xelatex
\documentclass[12pt,a4paper]{ctexart}
\usepackage{amsmath,amssymb,amsfonts}
\usepackage{graphicx}
\usepackage{booktabs}
\usepackage{hyperref}
\usepackage{geometry}
\usepackage{indentfirst}
\usepackage{caption}
\usepackage{enumitem}
\usepackage{listings}
\usepackage{xcolor}
\geometry{left=2.5cm,right=2.5cm,top=2.5cm,bottom=2.5cm}

% 正文小四号 ≈ 12pt；参考文献/附录五号
\newcommand{\bodyzi}{\zihao{-4}}
\newcommand{\refzi}{\zihao{5}}

\lstset{
  basicstyle=\ttfamily\small,
  breaklines=true,
  frame=single,
  backgroundcolor=\color{gray!8},
}

\begin{document}

% ========== 封面 ==========
\begin{titlepage}
  \centering
  \vspace*{1cm}
  {\Large \textbf{暨\hspace{1em}南\hspace{1em}大\hspace{1em}学}}\\[0.8cm]
  {\large 本科生课程论文}\\[2.5cm]
  {\LARGE \textbf{论文题目：}}\\[0.4cm]
  {\Large 2026 密码数学挑战赛赛题一：\\无穷范数短整数解（SIS$_{\infty}$）的\\格论启发式与整数规划混合求解研究}\\[2.5cm]
  \begin{tabular}{rl}
    学\hspace{2em}院： & \underline{\hspace{6cm}} \\[0.6cm]
    专\hspace{2em}业： & \underline{\hspace{6cm}} \\[0.6cm]
    学生姓名： & \underline{\hspace{6cm}} \\[0.6cm]
    学\hspace{2em}号： & \underline{\hspace{6cm}} \\[0.6cm]
    课程名称： & \underline{\hspace{6cm}} \\[0.6cm]
    指导教师： & \underline{\hspace{6cm}} \\
  \end{tabular}
  \vfill
  {\large \underline{\hspace{2cm}} 年 \underline{\hspace{1.5cm}} 月 \underline{\hspace{1.5cm}} 日}
\end{titlepage}

% ========== 摘要 ==========
\begin{abstract}
\noindent
2026 年全国高校密码数学挑战赛赛题一要求在模 $q$ 下求解短整数解：给定 $A\in\mathbb{Z}_q^{n\times m}$、$t\in\mathbb{Z}_q^n$，求 $u,v$ 使得 $Av+u\equiv t\pmod q$ 且 $\|u\|_{\infty},\|v\|_{\infty}\le\gamma$。官方将十道小问划分为\textbf{近似 SVP}、\textbf{近似 CVP}与\textbf{受限 SVP}三类；第三类另要求 $\|u\|_2^2+\|v\|_2^2<q^2$，禁止“先 $L_2$ 最短再过滤”\cite{wang2025}。

与 BKZ 2.0 的 $L_2$ 目标不同，赛题优化 Chebyshev 范数。本文提出\textbf{残差主导三阶段混合框架}：格种子自动调度（BKZ 2.0、G6K/BDGL 筛法、Kannan 嵌入、Wang enumerate-then-slice）$\to$ 字典序局部搜索 $\to$ CP-SAT 收尾，并对接阶梯计分\cite{challenge2026}。创新包括 Chebyshev 字典序 $(M,V,S)$、$O(n)$ 列增量更新、三类差异化流水线及 lex 欧氏收尾。标准全量实验 p1--p4 的 $E_\infty=42$--$44$。开源：\url{https://github.com/arriverl/sisinf}。

\vspace{0.5em}
\noindent\textbf{关键词：}短整数解；无穷范数；格基约化；受限最短向量；约束规划
\end{abstract}

\tableofcontents
\newpage

\bodyzi

% ==================== 1 引言 ====================
\section{引言}

短整数解（SIS）是格密码学核心假设\cite{ajtai1996,micciancio2009book}。赛题在
\begin{equation}\label{eq:main}
  Av+u\equiv t\pmod q,\qquad \|u\|_{\infty}\le\gamma,\quad \|v\|_{\infty}\le\gamma
\end{equation}
上附加无穷范数界，$n=m\in\{100,120,140,160\}$，$\gamma\in\{15,16,17,18\}$，属于\textbf{非平凡短向量区}。

难点有三：（1）BKZ/筛法优化 $L_2$，与 $L_{\infty}$ 目标失配；（2）十题分齐次/非齐次/欧氏受限三类；（3）官方\textbf{阶梯计分}按 $E_\infty=\max(\|u\|_{\infty},\|v\|_{\infty})$ 给 0--10 分\cite{challenge2026}，非二元可行判定。

本文对接 Chen--Nguyen BKZ 2.0\cite{chen2011}、Becker BDGL\cite{becker2016}、Kannan 嵌入\cite{kannan1987}与 Wang 受限 SVP\cite{wang2025}，形成可复现开源求解器。

% ==================== 2 原理或方案设计 ====================
\section{原理或方案设计}

\subsection{中心化与残差主导建模}

对称中心化 $\mathrm{center}:\mathbb{Z}^n\to(-q/2,q/2]$ 将模坐标映到最短代表元。定义
\begin{equation}\label{eq:residual}
  u=r(v)=\mathrm{center}(t-Av\bmod q).
\end{equation}
则式\eqref{eq:residual}自动满足同余；可行性等价于
\begin{equation}\label{eq:feas}
  \|r(v)\|_{\infty}\le\gamma \quad\land\quad \|v\|_{\infty}\le\gamma.
\end{equation}
\textbf{残差主导}将双变量搜索降为单变量 $v$，三类题共享。

\subsection{赛题三类划分}

\begin{table}[htbp]
  \centering
  \caption{赛题三类划分}
  \label{tab:taxonomy}
  \begin{tabular}{ccclc}
    \toprule
    类 & 小问 & $t$ & 官方路线 & 附加约束 \\
    \midrule
    一 & 1,3,6,9 & $\equiv 0$ & BKZ 2.0 + 筛法 & 非平凡 \\
    二 & 2,4,7,10 & $\not\equiv 0$ & Kannan / CVP & --- \\
    三 & 5,8 & 一般 & Wang 受限 SVP & $\|u\|_2^2+\|v\|_2^2<q^2$ \\
    \bottomrule
  \end{tabular}
\end{table}

\subsection{阶梯计分函数}

前提满足时（第三类另需 $L_2<q^2$）：
\begin{equation}\label{eq:score}
  S(E_\infty)=
  \begin{cases}
    10, & E_\infty\le\gamma;\\
    8,6,4,2, & E_\infty=\gamma{+}1,{\ldots},\gamma{+}4;\\
    0, & E_\infty>\gamma{+}4.
  \end{cases}
\end{equation}

\subsection{格论背景}

\paragraph{BKZ 2.0与BDGL筛法}
BKZ 以块大小 $\beta$ 约化 Ajtai 格基\cite{chen2011}；BDGL 筛法\cite{becker2016}与 BKZ 组合为当前实用最强近似 SVP 路线。本文经 G6K 接入 \texttt{bdgl2}，无 G6K 时回退 list sieve。

\paragraph{Kannan嵌入（第二类）}
将 CVP 嵌入 $(n{+}m{+}1)$ 维 SVP，搜索接近 $t$ 的短向量后模拉回\cite{kannan1987}。

\paragraph{Wang受限SVP（第三类）}
\textbf{enumerate-then-slice}：尾块枚举生成短向量列表 $\to$ 按 $L_{\infty}$ 与 $<q^2$ slice $\to$ dimension-for-free 降维\cite{wang2025}。禁止 $L_2$ 最短再过滤。

\subsection{三阶段混合架构}

\begin{enumerate}[leftmargin=2em]
  \item \textbf{格种子}：BKZ / G6K / sieve / Kannan / Wang slice；
  \item \textbf{字典序局部搜索}：多重重启 + 邻域算子；
  \item \textbf{CP-SAT收尾}：\texttt{full}（类一、二）或 \texttt{lex}（类三）。
\end{enumerate}

调度：\texttt{apply\_sis\_class\_defaults}；全量：\texttt{--full-max}。

\begin{table}[htbp]
  \centering
  \caption{三类流水线差异}
  \label{tab:pipeline}
  \small
  \begin{tabular}{p{1.4cm}p{3.8cm}p{3.8cm}p{3.6cm}}
    \toprule
    阶段 & 第一类 & 第二类 & 第三类 \\
    \midrule
    格种子 & BKZ+G6K/sieve+Wang & Kannan+CVP & G6K+Wang slice \\
    搜索 & 核游走、$u$ snap & pull kick & 稀疏化、$L_2$ 惩罚 \\
    收尾 & full CP-SAT & full CP-SAT & lex+抛光 \\
    \bottomrule
  \end{tabular}
\end{table}

\subsection{Chebyshev字典序目标}

溢出 $o^u_i=\max(|u_i|-\gamma,0)$，$o^v_j=\max(|v_j|-\gamma,0)$。定义
\begin{equation}\label{eq:lex}
  V=\|o^u\|_0+\|o^v\|_0,\quad
  S=\|o^u\|_1+\|o^v\|_1,\quad
  M=\max(\|o^u\|_{\infty},\|o^v\|_{\infty}).
\end{equation}
字典序：$(M,V,S)$，与阶梯计分对齐。

搜索能量（连续代理）：
\begin{equation}\label{eq:energy}
  \mathcal{E}=10^6 V + w_M M + w_S S + w_e e_{\mathrm{excess}} - w_H H,
\end{equation}
其中 $e_{\mathrm{excess}}=\max(0,\|u\|_2^2+\|v\|_2^2-q^2{+}1)$（第三类），$H$ 为分箱熵。

\subsection{$O(n)$列增量更新}

修改 $v_j\leftarrow v_j+\delta$ 时：
\begin{equation}\label{eq:delta}
  r'(v)=\mathrm{center}\bigl(r(v)-\delta\cdot A_{\cdot j}\bigr),
\end{equation}
邻域试探 $O(n)$，支撑百万级搜索。

\subsection{三类差异化方案}

\paragraph{第一类（1,3,6,9）}
$\text{BKZ}(\beta{\approx}28\text{--}56)\to\text{G6K/sieve}\to\text{Wang }L_{\infty}\text{ slice}\to\text{Wagner}\to\text{核游走}\to\text{full CP-SAT}$。

\paragraph{第二类（2,4,7,10）}
Kannan 嵌入 + BKZ；CVP 提升与模拉回；关闭 BKZ 主种子；pull kick $\Delta v\propto -A^\top\mathrm{sign}(u)$。

\paragraph{第三类（5,8）}
短向量列表 $\to$ slice($L_{\infty}$,$<q^2$) $\to$ lex CP-SAT $\to$ \texttt{euclid\_polish}。

\subsection{CP-SAT建模}

$v_j\in[-\gamma,\gamma]$，slack $k_i$ 线性化 $(Av)_i+u_i-qk_i=t_i$。模式：\texttt{full}、\texttt{chunk}、\texttt{lex}。

\subsection{创新点归纳}

\begin{enumerate}[leftmargin=2em]
  \item 残差主导统一建模（式\eqref{eq:residual}）；
  \item Chebyshev 字典序（式\eqref{eq:lex}）与阶梯计分（式\eqref{eq:score}）对齐；
  \item 三类自动调度与 \texttt{--full-max} 全量预设；
  \item 格种子分层栈（BKZ/G6K/sieve/Kannan/Wang）；
  \item $O(n)$ 列增量（式\eqref{eq:delta}）；
  \item 第三类 lex 双约束收尾；
  \item 文献最优—实现—可达三层分解。
\end{enumerate}

% ==================== 3 程序实现或算法分析 ====================
\section{程序实现或算法分析}

\subsection{软件架构}

核心模块见表~\ref{tab:modules}。主链：\texttt{run\_full\_validation} $\to$ \texttt{local\_search\_one} $\to$ \texttt{execute\_finish} $\to$ \texttt{verify\_solution}。

\begin{table}[htbp]
  \centering
  \caption{核心模块}
  \label{tab:modules}
  \small
  \begin{tabular}{p{3.6cm}p{4.4cm}p{4.8cm}}
    \toprule
    模块 & 功能 & 类别 \\
    \midrule
    \texttt{solve\_sisinf.py} & 主求解器、SearchConfig & 全部 \\
    \texttt{lattice\_bkz} & LLL/BKZ 2.0 & 一、三 \\
    \texttt{lattice\_g6k} & G6K BDGL2 & \texttt{--full-max} \\
    \texttt{lattice\_sieve} & list sieve & 一、三 \\
    \texttt{lattice\_kannan} & Kannan 嵌入 & 二 \\
    \texttt{lattice\_restricted\_svp} & Wang enumerate-slice & 一、三 \\
    \texttt{sis\_scoring} & 阶梯计分 & 验证 \\
    \texttt{run\_full\_validation.py} & 十题评测 & 全部 \\
    \bottomrule
  \end{tabular}
\end{table}

\subsection{格种子实现要点}

\textbf{BKZ}：构造 $(n{+}m)$ 维 Ajtai 基，fpylll 多轮 BKZ，从约化列提取 $v$。

\textbf{G6K}：投影尾块执行 \texttt{bdgl2}，饱和检测与 BKZ 回退。

\textbf{Wang slice}：\texttt{enumerate\_approx\_svp\_list} $\to$ \texttt{slice\_by\_linf\_and\_norm}；参数 \texttt{wang\_enum\_tail\_rank}、\texttt{pool\_size}、\texttt{max\_trials} 控制算力。

\subsection{局部搜索算法}

\textbf{算法1}（单 restart 概要）：初始化 $r=\mathrm{center}(t{-}Av_0)$；每步随机排列坐标 $j$；对 $\delta\in[-\Delta,\Delta]$ 用式\eqref{eq:delta} 试探；字典序接受；停滞时 pull kick / 核游走 / 块 CP。

邻域算子：单坐标贪心、pull kick、核游走、pair relief、$u$ 行 snap、block CP。

\subsection{复杂度与算力}

\begin{table}[htbp]
  \centering
  \caption{复杂度概估}
  \label{tab:complexity}
  \small
  \begin{tabular}{lcc}
    \toprule
    组件 & 复杂度 & 依赖 \\
    \midrule
    BKZ 2.0 & $O(n^3\cdot 2^{O(\beta)})$ 启发式 & fpylll \\
    G6K bdgl2 & 指数（维数相关） & Cython, 多核 \\
    Wang enumerate & $O(\binom{r}{k}(2c{+}1)^k)$ & 尾块秩 $r$ \\
    局部搜索 & $O(Tm\Delta n)$ & numpy \\
    full CP-SAT & 最坏指数 & OR-Tools \\
    \bottomrule
  \end{tabular}
\end{table}

推荐：Linux，16+ 核，32GB+ RAM；\texttt{--full-max} 建议 64GB。

\subsection{环境安装}

\begin{lstlisting}[language=bash]
cd /home/cwh/sisinf
conda activate cwh
conda install -c conda-forge fpylll -y
pip install numpy ortools "Cython>=3.0" cysignals
bash scripts/install_g6k.sh
python3 scripts/sis_cli.py check all
python3 scripts/sis_cli.py check g6k
\end{lstlisting}

\subsection{运行命令}

\textbf{（1）冒烟：}\texttt{python3 scripts/sis\_cli.py check smoke}

\textbf{（2）快速十题：}\texttt{sis\_cli.py --quick}（1--3\,h）

\textbf{（3）标准全量（推荐）：}
\begin{lstlisting}[language=bash]
python3 scripts/sis_cli.py \
  --batch-rounds 6 --ilp-time-limit 3600
\end{lstlisting}

\textbf{（4）论文全量：}\texttt{--full-max --batch-rounds 24 --ilp-time-limit 14400}

\textbf{（5）分三类：}\texttt{run\_class\_batch.py --class \{1,2,3\} --full-max}

\textbf{（6）单题：}\texttt{--problems 1 --batch-rounds 6 --ilp-time-limit 3600}

报告：\texttt{results/full\_validation/full\_validation\_report.json}。

\subsection{实验环境与配置}

Linux 服务器（conda \texttt{cwh}）；fpylll + ortools；标准全量（非 \texttt{--full-max}）：
\texttt{--batch-rounds 6 --ilp-time-limit 3600}。

\subsection{十题全量验证结果}

\begin{table}[htbp]
  \centering
  \caption{十题标准全量验证（p1--p4 已完成，p5--p10 进行中）}
  \label{tab:full10}
  \small
  \begin{tabular}{cccccccccc}
    \toprule
    题号 & 类 & $\gamma$ & $n$ & 启发式 & ILP后 & $\|v\|_\infty$ & 同余 & norm & 得分 & 状态 \\
    \midrule
    1 & 一 & 15 & 100 & 44 & \textbf{42} & 15 & $\checkmark$ & $\checkmark$ & 0 & 完成 \\
    2 & 二 & 15 & 100 & 42 & \textbf{42} & 15 & $\checkmark$ & $\checkmark$ & 0 & 完成 \\
    3 & 一 & 16 & 120 & 45 & \textbf{43} & 15 & $\checkmark$ & $\checkmark$ & 0 & 完成 \\
    4 & 二 & 16 & 120 & 45 & \textbf{44} & 15 & $\checkmark$ & $\checkmark$ & 0 & 完成 \\
    5 & 三 & 16 & 100 & 45 & --- & 8--12 & $\checkmark$ & $\times$ & --- & batch \\
    6--10 & --- & --- & --- & --- & --- & --- & --- & --- & --- & 待完成 \\
    \bottomrule
  \end{tabular}
\end{table}

\noindent
\textbf{ILP 降幅：}题1 $44{\to}42$；题2 $42{\to}42$；题3 $45{\to}43$；题4 $45{\to}44$。
距可行 $E_\infty{-}\gamma$：27、27、27、28；得分均为 0。

\begin{table}[htbp]
  \centering
  \caption{第一类（题 1、3）明细}
  \label{tab:class1}
  \begin{tabular}{ccccccc}
    \toprule
    题号 & $q$ & 启发式 & ILP后 & $\|u\|_\infty$ & $\|v\|_\infty$ & 得分 \\
    \midrule
    1 & 100 & 44 & \textbf{42} & 42 & 15 & 0 \\
    3 & 120 & 45 & \textbf{43} & 43 & 15 & 0 \\
    \bottomrule
  \end{tabular}
\end{table}

\begin{table}[htbp]
  \centering
  \caption{第二类（题 2、4）明细}
  \label{tab:class2}
  \begin{tabular}{ccccccc}
    \toprule
    题号 & $q$ & 启发式 & ILP后 & $\|u\|_\infty$ & $\|v\|_\infty$ & 得分 \\
    \midrule
    2 & 100 & 42 & \textbf{42} & 42 & 15 & 0 \\
    4 & 120 & 45 & \textbf{44} & 44 & 15 & 0 \\
    \bottomrule
  \end{tabular}
\end{table}

同余满足；瓶颈在 $\|u\|_{\infty}$；ILP 1\,h 后平台；题5 batch 阶段 norm\_sq $\approx 7.3{\times}10^4 \gg q^2$。

\begin{table}[htbp]
  \centering
  \caption{可行可得分解（标准全量，p1--p4）}
  \label{tab:decomp}
  \small
  \begin{tabular}{p{1cm}p{3cm}p{3.6cm}p{3.8cm}}
    \toprule
    类 & 文献最优 & 本文实现 & 当前可达 \\
    \midrule
    一 & BKZ+BDGL & BKZ+sieve+Wang+CP-SAT & $E_\infty{=}42$--$43$ \\
    二 & Kannan+BKZ & Kannan+CVP+CP-SAT & $E_\infty{=}42$--$44$ \\
    三 & Wang 受限 SVP & slice+lex+polish & 题5 batch，norm 未达标 \\
    \bottomrule
  \end{tabular}
\end{table}

% ==================== 4 总结 ====================
\section{总结}

本文针对 SIS$_{\infty}$ 赛题，在残差主导建模下提出格种子、字典序搜索与 CP-SAT 收尾的三阶段混合框架\cite{sisinfrepo}。贡献：统一式\eqref{eq:residual}；字典序式\eqref{eq:lex} 对齐阶梯计分；工程化 BKZ/G6K/Kannan/Wang 栈；$O(n)$ 增量式\eqref{eq:delta}；开源可复现。

实验：标准全量下 p1--p4 的 $E_\infty$ 为 42--44；全维 CP-SAT 为主要压降；十题报告完成后补全表~\ref{tab:full10}。

\textbf{展望：}（1）十题 \texttt{--full-max} 报告；（2）G6K 调参；（3）lex 欧氏硬约束；（4）与 BKZ 基线对比。

% ==================== 参考文献（五号） ====================
\refzi
\begin{thebibliography}{99}

\bibitem{ajtai1996}
M. Ajtai. Generating hard instances of lattice problems. \emph{STOC}, 1996.

\bibitem{lenstra1982}
A. K. Lenstra, H. W. Lenstra Jr., L. Lov\'asz. Factoring polynomials with rational coefficients. \emph{Math. Ann.}, 261:515--534, 1982.

\bibitem{micciancio2009book}
D. Micciancio, O. Regev. Lattice-based cryptography. In \emph{Post-Quantum Cryptography}, Springer, 2009.

\bibitem{chen2011}
Y. Chen, P. Q. Nguyen. BKZ 2.0: Better lattice security estimates. \emph{ASIACRYPT}, 2011.

\bibitem{becker2016}
A. Becker, L. Ducas, N. Gama, M. Tibouchi. New directions in nearest neighbor searching with applications to lattice sieving. \emph{SODA}, 2016.

\bibitem{wagner2002}
D. Wagner. A generalized birthday problem. \emph{CRYPTO}, 2002.

\bibitem{kannan1987}
R. Kannan. Minkowski's convex body theorem and integer programming. \emph{Math. Oper. Res.}, 12(3):415--440, 1987.

\bibitem{babai1986}
L. Babai. On Lov\'asz's lattice reduction. \emph{Combinatorica}, 6(1):1--13, 1986.

\bibitem{achterberg2007}
T. Achterberg. Constraint Integer Programming. PhD thesis, TU Berlin, 2007.

\bibitem{ortools2024}
Google OR-Tools. CP-SAT documentation, 2024.

\bibitem{fpylll2019}
The fpylll team. fpylll, 2019.

\bibitem{wang2025}
G. Wang, W. Xia, D. Gu. Heuristic algorithm for restricted SVP. \emph{PQCrypto}, 2025. \url{https://eprint.iacr.org/2025/586}

\bibitem{challenge2026}
2026 全国高校密码数学挑战赛. 赛题一解析与评分规则, 2026.

\bibitem{sisinfrepo}
arriverl. \emph{sisinf}. GitHub, 2026. \url{https://github.com/arriverl/sisinf}

\bibitem{hoos2004}
H. H. Hoos, T. St\"utzle. \emph{Stochastic Local Search}. Morgan Kaufmann, 2004.

\end{thebibliography}

% ==================== 附录（五号） ====================
\appendix
\section{十题参数一览}
\refzi
\begin{table}[htbp]
  \centering
  \small
  \begin{tabular}{cccccc}
    \toprule
    题号 & 类 & $n{=}m$ & $q$ & $\gamma$ & 欧氏约束 \\
    \midrule
    1 & 一 & 100 & 100 & 15 & 否 \\
    2 & 二 & 100 & 100 & 15 & 否 \\
    3 & 一 & 120 & 120 & 16 & 否 \\
    4 & 二 & 120 & 120 & 16 & 否 \\
    5 & 三 & 100 & 100 & 16 & 是 \\
    6 & 一 & 140 & 140 & 17 & 否 \\
    7 & 二 & 140 & 140 & 17 & 否 \\
    8 & 三 & 120 & 120 & 18 & 是 \\
    9 & 一 & 160 & 160 & 18 & 否 \\
    10 & 二 & 160 & 160 & 18 & 否 \\
    \bottomrule
  \end{tabular}
\end{table}

\section{SearchConfig 三类默认差异}
\refzi
\begin{table}[htbp]
  \centering
  \small
  \begin{tabular}{lccc}
    \toprule
    参数 & 类一 & 类二 & 类三 \\
    \midrule
    \texttt{use\_bkz\_seeds} & $\checkmark$ & $\times$ & $\checkmark$ \\
    \texttt{use\_kannan\_seeds} & $\times$ & $\checkmark$ & $\times$ \\
    \texttt{use\_restricted\_svp} & $\checkmark$ & $\times$ & $\checkmark$ \\
    \texttt{bkz\_beta} & 28--32 & 24--28 & 24--28 \\
    \texttt{euclid\_weight} & 0.5 & 0.6 & 3--4 \\
    CP-SAT & full & full & lex \\
    \bottomrule
  \end{tabular}
\end{table}

\section{阶梯计分伪代码}
\refzi
\begin{lstlisting}[language=Python]
def competition_score(gamma, inf_u, inf_v, congruence_ok,
                      norm_sq, q, sis_class):
    if not congruence_ok: return 0
    if sis_class == 3 and norm_sq >= q * q: return 0
    e_inf = max(inf_u, inf_v)
    if e_inf <= gamma: return 10
    if e_inf == gamma + 1: return 8
    if e_inf == gamma + 2: return 6
    if e_inf == gamma + 3: return 4
    if e_inf == gamma + 4: return 2
    return 0
\end{lstlisting}

\section{LaTeX 编译}
\refzi
\begin{verbatim}
cd docs && xelatex COURSE_PAPER.tex && xelatex COURSE_PAPER.tex
\end{verbatim}

\end{document}
